% META: {"id": "TEXP-ARI-EX-031", "chapitre": "TEXP-ARITHMETIQUE", "type_objet": "exercice", "capacites": ["TEXP-ARI-C4"], "capacites_codes": ["C4"], "methodes": ["M4"], "parcours": 2, "competences": ["raisonner","calculer"], "duree_min": 15, "mode_creation": "ex_nihilo", "sources_inspiration": [], "fichier_tex": "chapitres/TEXP-ARITHMETIQUE/exercices/TEXP-ARI-EX-031.tex", "corrige_tex": "chapitres/TEXP-ARITHMETIQUE/corriges/TEXP-ARI-CO-031.tex", "status": "approved"}
% BEGIN-VERIFY
% from sympy import *
% x = symbols('x')
% # f''(x) = 6(x-1)(x-3)
% fpp = 6*(x-1)*(x-3)
% # f' = x^3 - 6x^2 + 9x + C, f = x^4/4 - 2x^3 + 9x^2/2 + Cx + D
% fp = integrate(fpp, x)
% assert simplify(fp - (2*x**3 - 12*x**2 + 18*x)) == 0
% f = integrate(fp, x)
% assert simplify(f - (x**4/2 - 4*x**3 + 9*x**2)) == 0
% assert solve(fpp, x) == [1, 3]
% END-VERIFY
\begin{exercice}{TEXP-ARI-EX-031}{2}{15}
On sait que $f''(x) = 6(x-1)(x-3)$ et que $f(0) = 0$, $f'(0) = 0$.
\begin{enumerate}
  \item Determiner une expression de $f$.
  \item Identifier les intervalles de convexite et les points d'inflexion.
  \item Esquisser la courbe de $f$.
\end{enumerate}

\end{exercice}
